\documentclass[11pt,a4paper]{article}
% preamble.tex — shared setup for human-verification documents
% Compile each document with:  tectonic <File>.tex

\usepackage[T1]{fontenc}
\usepackage[utf8]{inputenc}
\usepackage{lmodern}
\usepackage{amssymb,amsmath,amsthm}
\usepackage{mathtools}
\usepackage[dvipsnames]{xcolor}
\usepackage{listings}
\usepackage{hyperref}
\usepackage[margin=2.5cm]{geometry}
\usepackage{parskip}

% ── Colours ──────────────────────────────────────────────────────────
\definecolor{leanBg}{HTML}{F7F7F0}
\definecolor{haskellBg}{HTML}{F0F4F7}
\definecolor{keyword}{HTML}{7B3294}
\definecolor{comment}{HTML}{5A8C32}
\definecolor{noteFg}{HTML}{B35900}

% ── Lean 4 listings language ─────────────────────────────────────────
% NOTE: Lean listing blocks use ASCII approximations of unicode symbols
% (e.g. -> for arrow, o. for compose).  The literate key substitutes
% these with proper math-mode renderings in the PDF.
\lstdefinelanguage{Lean4}{
  morekeywords={def,theorem,structure,where,namespace,end,import,fun,
                let,have,show,by,exact,intro,cases,rfl,subst,obtain,
                funext,simp,rw,apply,variable,Type,Prop,instance,
                section,open,class,extends,deriving,sorry,noncomputable},
  sensitive=true,
  morecomment=[l]{--},
  morecomment=[n]{/-}{-/},
  morestring=[b]",
}

% ── Lean listing style ───────────────────────────────────────────────
\lstdefinestyle{lean}{
  language=Lean4,
  backgroundcolor=\color{leanBg},
  basicstyle=\small\ttfamily,
  keywordstyle=\color{keyword}\bfseries,
  commentstyle=\color{comment}\itshape,
  breaklines=true,
  frame=single,
  framerule=0.4pt,
  xleftmargin=4pt,
  xrightmargin=4pt,
  aboveskip=6pt,
  belowskip=4pt,
  literate=
    {->}{{$\to$}}2
    {=>}{{$\Rightarrow$}}2
    {<|}{{$\langle$}}1
    {|>}{{$\rangle$}}1
    {==>}{{$\Longrightarrow$}}3
    { o. }{{$\;\circ\;$}}3
    {forall}{{$\forall$}}1
    {:=}{{$\,:=\,$}}2
    {!=}{{$\neq$}}2
    {<=}{{$\leq$}}2
    {>=}{{$\geq$}}2,
}

% ── Haskell listing style ───────────────────────────────────────────
\lstdefinestyle{haskell}{
  language=Haskell,
  backgroundcolor=\color{haskellBg},
  basicstyle=\small\ttfamily,
  keywordstyle=\color{keyword}\bfseries,
  commentstyle=\color{comment}\itshape,
  breaklines=true,
  frame=single,
  framerule=0.4pt,
  xleftmargin=4pt,
  xrightmargin=4pt,
  aboveskip=6pt,
  belowskip=4pt,
  literate=
    {->}{{$\to$}}2
    {=>}{{$\Rightarrow$}}2,
}

% ── Theorem environments ────────────────────────────────────────────
\theoremstyle{definition}
\newtheorem{defn}{Definition}[section]
\newtheorem{thm}[defn]{Theorem}
\newtheorem{prop}[defn]{Proposition}

\theoremstyle{remark}
\newtheorem*{fidelitynote}{Fidelity note}

% ── Convenience macros ───────────────────────────────────────────────
\newcommand{\lean}[1]{\texttt{#1}}      % inline Lean identifier
\newcommand{\hask}[1]{\texttt{#1}}      % inline Haskell identifier
\newcommand{\Ext}{\mathrm{Ext}}
\newcommand{\Pos}{\mathrm{Pos}}
\newcommand{\Shape}{\mathrm{Shape}}


\title{Human Verification: \texttt{TrustBoundary}\\[4pt]
  \large Trust-Labelled Trees, the Boundary Rule, and Duplicate-Stability}
\author{Machine-generated from \texttt{lean/TrustBoundary.lean} for review by Neil Ghani}
\date{July 2026}

\begin{document}
\maketitle

\begin{abstract}
This document presents every definition and theorem \emph{statement} from the
Lean~4 module \lean{TrustBoundary} in three parallel forms: the original
Lean, a readable Haskell rendering, and precise mathematical English.  Proofs
are omitted --- they are machine-checked by Lean.

The module links proof registries (finitely-branching rose trees with trust
labels) to the trust-system theory of \lean{Trust.lean}.  The main structural
result is \emph{duplicate-stability}: if the global boundary rule holds for
a tree $t$, it also holds for every subtree $t.\mathrm{sub}(p)$.  The
deployed trust chain (\lean{TrustLevel}, 7 levels) and the loop experiment's
chain (\lean{LoopTrustLevel}, 6 levels) are also instantiated here.
\end{abstract}

\tableofcontents
\bigskip

%===================================================================
\section{Overview}
%===================================================================

Deployed proof registries are finitely-branching rose trees (\lean{LTree})
whose nodes carry a trust label.  The boundary rule
(\lean{registry\_validate.py}) is node-local: a node may claim
\lean{proved} or above only if every non-dead child claims \lean{proved} or
above.  Equivalently, the rule is a \emph{coKleisli predicate}: apply the
local check at every node produced by the comonad operation \lean{duplicate}
(= \lean{sub}).

The module proves:
\begin{enumerate}
  \item \textbf{Duplicate-stability} (\lean{ok\_sub}): the global boundary
    rule propagates to every subtree.
  \item \textbf{The deployed chain} (\lean{TrustLevel}): a \lean{Chain}
    instance for the seven-level registry chain, so the certificate lemma
    and lax transport from \lean{Trust.lean} apply to real registry labels.
  \item \textbf{The loop experiment's chain} (\lean{LoopTrustLevel}): a
    six-level variant for the Warnaar positivity loop.
\end{enumerate}

%===================================================================
\section{Trust labels and labelled trees}
%===================================================================

\subsection{Structure: \lean{TrustLabel}}

\begin{defn}[Trust label]
A \textbf{trust label} is a pair: a trust level $\ell \in T$ (a level in
the chain $T$) and a boolean \lean{dead} flag.  Dead-end nodes are
\emph{exempt} from the boundary rule; \lean{dead = true} means the approach
was abandoned.  An order relation cannot express an exemption, which is why
\lean{dead} is a separate flag rather than a chain element.
\end{defn}

\paragraph{Lean.}
\begin{lstlisting}[style=lean]
structure TrustLabel (T : Type) where
  trust : T
  dead  : Bool
\end{lstlisting}

\paragraph{Haskell.}
\begin{lstlisting}[style=haskell]
data TrustLabel t = TrustLabel
  { trust :: t
  , dead  :: Bool
  }
\end{lstlisting}


\subsection{Inductive type: \lean{LTree}}

\begin{defn}[Labelled rose tree]
A \textbf{trust-labelled finitely-branching rose tree} has a
\lean{TrustLabel T} at every node and an explicit arity $n : \mathbb{N}$
with children indexed by $\mathrm{Fin}(n)$.  This is the shape of a deployed
registry (\texttt{proofs/registry/*.json}).
\end{defn}

\paragraph{Lean.}
\begin{lstlisting}[style=lean]
inductive LTree (T : Type) where
  | node : TrustLabel T -> (n : Nat) -> (Fin n -> LTree T) -> LTree T
\end{lstlisting}

\paragraph{Haskell.}
\begin{lstlisting}[style=haskell]
data LTree t = Node (TrustLabel t) [LTree t]
-- (Using a list of children rather than a Fin-indexed function)
\end{lstlisting}

\begin{fidelitynote}
In Lean the arity $n$ is an explicit \lean{Nat} and children are indexed by
\lean{Fin n}, giving a \emph{dependent} function type.  In Haskell we use a
list, losing the arity-indexed dependency.
\end{fidelitynote}


\subsection{Definition: \lean{LTree.lab}}

\begin{defn}[Root label]
$\mathrm{lab}(t)$ extracts the trust label at the root of $t$.  Under the
comonad, $\mathrm{lab}(t.\mathrm{sub}(p))$ reads the label at position~$p$.
\end{defn}

\paragraph{Lean.}
\begin{lstlisting}[style=lean]
def LTree.lab : LTree T -> TrustLabel T
  | .node l _ _ => l
\end{lstlisting}

\paragraph{Haskell.}
\begin{lstlisting}[style=haskell]
lab :: LTree t -> TrustLabel t
lab (Node l _) = l
\end{lstlisting}

%===================================================================
\section{Positions, sub-trees, and paths}
%===================================================================

\subsection{Inductive type: \lean{LTree.LPos}}

\begin{defn}[Position / path]
A \textbf{position} in $t : \lean{LTree}\,T$ is a path from the root to a
node.  It is either \lean{here} (the root) or a child index $i$ followed by
a position in the $i$-th child.
\end{defn}

\paragraph{Lean.}
\begin{lstlisting}[style=lean]
inductive LPos : LTree T -> Type where
  | here  : LPos t
  | child : (i : Fin n) -> LPos (children i)
              -> LPos (.node l n children)
\end{lstlisting}

\paragraph{Haskell.}
\begin{lstlisting}[style=haskell]
data LPos t
  = Here
  | Child Int (LPos t)   -- approximate: index is Int, not Fin n
\end{lstlisting}


\subsection{Definition: \lean{LTree.sub}}

\begin{defn}[Sub-tree]
$t.\mathrm{sub}(p)$ is the sub-tree rooted at position $p$.
\end{defn}

\paragraph{Lean.}
\begin{lstlisting}[style=lean]
def sub : (t : LTree T) -> LPos t -> LTree T
  | t, .here => t
  | .node _ _ children, .child i p => (children i).sub p
\end{lstlisting}

\paragraph{Haskell.}
\begin{lstlisting}[style=haskell]
sub :: LTree t -> LPos (LTree t) -> LTree t
sub t Here           = t
sub (Node _ cs) (Child i p) = sub (cs !! i) p
\end{lstlisting}


\subsection{Definition: \lean{LTree.shift}}

\begin{defn}[Path concatenation]
$t.\mathrm{shift}(p, q)$ concatenates path $p$ with path $q$ (where $q$ is
a position in $t.\mathrm{sub}(p)$), giving a position in $t$.
\end{defn}

\paragraph{Lean.}
\begin{lstlisting}[style=lean]
def shift : (t : LTree T) -> (p : LPos t)
              -> LPos (t.sub p) -> LPos t
  | _, .here, q => q
  | .node _ _ children, .child i p, q =>
      .child i ((children i).shift p q)
\end{lstlisting}

\paragraph{Haskell.}
\begin{lstlisting}[style=haskell]
shift :: LTree t -> LPos (LTree t) -> LPos (LTree t) -> LPos (LTree t)
shift _ Here q = q
shift (Node _ cs) (Child i p) q = Child i (shift (cs !! i) p q)
\end{lstlisting}

%===================================================================
\section{Directed-container axioms for \lean{LTree}}
%===================================================================

\subsection{Definitions: \lean{shift\_root}, \lean{sub\_shift}, \lean{shift\_assoc}}

The following three lemmas confirm that \lean{LTree} with \lean{sub} and
\lean{shift} satisfies the directed-container axioms D3, D4, D5.

\begin{thm}[D3: shift to root]
For every tree $t$ and position $p$,
\[
  t.\mathrm{shift}(p, \mathrm{here}) = p.
\]
\end{thm}

\paragraph{Lean.}
\begin{lstlisting}[style=lean]
def shift_root : (t : LTree T) -> (p : LPos t) ->
    t.shift p .here = p
\end{lstlisting}

\begin{thm}[D4: sub-trees nest along shift]
For every $t$, $p : \mathrm{LPos}(t)$, $q : \mathrm{LPos}(t.\mathrm{sub}(p))$,
\[
  t.\mathrm{sub}(t.\mathrm{shift}(p, q)) = (t.\mathrm{sub}(p)).\mathrm{sub}(q).
\]
\end{thm}

\paragraph{Lean.}
\begin{lstlisting}[style=lean]
def sub_shift : (t : LTree T) -> (p : LPos t)
    -> (q : LPos (t.sub p)) ->
    t.sub (t.shift p q) = (t.sub p).sub q
\end{lstlisting}

\begin{thm}[D5: shift is associative]
For $t$, $p$, $q$, $r$ at compatible types (with transport along D4),
\[
  t.\mathrm{shift}(t.\mathrm{shift}(p, q),\, r) =
  t.\mathrm{shift}(p,\, (t.\mathrm{sub}(p)).\mathrm{shift}(q, r)).
\]
\end{thm}

\paragraph{Lean.}
\begin{lstlisting}[style=lean]
def shift_assoc : (t : LTree T) -> (p : LPos t)
    -> (q : LPos (t.sub p))
    -> (r : LPos ((t.sub p).sub q)) ->
    t.shift (t.shift p q) ((t.sub_shift p q).symm => r) =
      t.shift p ((t.sub p).shift q r)
\end{lstlisting}

\begin{fidelitynote}
The \lean{=>} in D5 is Lean transport notation: $(t.\mathrm{sub\_shift}\,p\,q)^{-1}
\blacktriangleright r$ re-types $r : \mathrm{LPos}((t.\mathrm{sub}(p)).\mathrm{sub}(q))$
as $r : \mathrm{LPos}(t.\mathrm{sub}(t.\mathrm{shift}(p,q)))$ via the equality from D4.
In Haskell this coercion is invisible (all trees have the same type), but
in Lean's dependent type theory the transport is an explicit proof term.
The listings use \texttt{=>} as an ASCII approximation of the transport operator.
\end{fidelitynote}

\paragraph{Haskell.}
\begin{lstlisting}[style=haskell]
-- shift_root:  shift t p Here = p
-- sub_shift:   sub t (shift t p q) = sub (sub t p) q
-- shift_assoc: shift t (shift t p q) r = shift t p (shift (sub t p) q r)
-- All hold by structural induction on p.
\end{lstlisting}

%===================================================================
\section{The trust registry as a directed container}
%===================================================================

\subsection{Definition: \lean{trustRegistry}}

\begin{defn}[Trust registry directed container]
The trust-labelled rose trees form a directed container with
\begin{align*}
  \mathrm{Shape} &= \lean{LTree}\,T, \\
  \mathrm{Pos}   &= \lean{LTree.LPos}, \\
  \mathrm{root}  &= \lambda t.\,\lean{here}, \\
  \mathrm{sub}   &= \lean{LTree.sub}, \\
  \mathrm{shift} &= \lean{LTree.shift}.
\end{align*}
Axioms D1--D5 are provided by \lean{rfl}, \lean{sub\_shift}, \lean{shift\_root},
and \lean{shift\_assoc} respectively (D1 and D2 hold definitionally).
\end{defn}

\paragraph{Lean.}
\begin{lstlisting}[style=lean]
def trustRegistry (T : Type) : DirectedContainer where
  Shape := LTree T
  Pos   := LTree.LPos
  root  := fun _ => .here
  sub   := LTree.sub
  shift := LTree.shift
  d1    := fun _ => rfl
  d4    := LTree.sub_shift
  d2    := fun _ _ => rfl
  d3    := LTree.shift_root
  d5    := LTree.shift_assoc
\end{lstlisting}

\paragraph{Haskell.}
\begin{lstlisting}[style=haskell]
-- trustRegistry :: DirectedContainer (LTree T) (LPos (LTree T))
-- root t = Here
-- sub  = LTree.sub
-- shift = LTree.shift
-- D1--D5 hold by structural induction.
\end{lstlisting}

%===================================================================
\section{The boundary rule}
%===================================================================

\subsection{Definition: \lean{LTree.localOk}}

\begin{defn}[Node-local boundary rule]
At threshold $\theta : T$, a tree $t = \lean{node}\,l\,n\,\mathrm{children}$
satisfies the \textbf{local boundary rule} if: whenever $\theta \le l.\mathrm{trust}$
(this node claims above threshold), every non-dead child $i$ satisfies
$\theta \le \mathrm{children}(i).\mathrm{lab}.\mathrm{trust}$.

Formally:
\[
  \mathrm{localOk}(\theta, t) \;\iff\;
  \theta \le t.\mathrm{lab}.\mathrm{trust} \;\Rightarrow\;
  \forall\, i : \mathrm{Fin}(n),\;
  \mathrm{children}(i).\mathrm{lab}.\mathrm{dead} = \mathrm{true}
  \;\vee\; \theta \le \mathrm{children}(i).\mathrm{lab}.\mathrm{trust}.
\]
This is precisely what \texttt{registry\_validate.py} checks at each node.
\end{defn}

\paragraph{Lean.}
\begin{lstlisting}[style=lean]
def localOk (theta : T) : LTree T -> Prop
  | .node l n children =>
      theta <= l.trust -> forall i : Fin n,
        (children i).lab.dead = true \/ theta <= (children i).lab.trust
\end{lstlisting}

\paragraph{Haskell.}
\begin{lstlisting}[style=haskell]
localOk :: (Ord t) => t -> LTree t -> Bool
localOk theta (Node l children)
  | trust l >= theta = all (\c -> dead (lab c) || trust (lab c) >= theta) children
  | otherwise        = True
\end{lstlisting}


\subsection{Definition: \lean{LTree.Ok}}

\begin{defn}[Global boundary rule]
The \textbf{global boundary rule} holds for $t$ at threshold $\theta$ if the
local check holds at every position (i.e., at every sub-tree produced by
\lean{duplicate}):
\[
  \mathrm{Ok}(\theta, t) \;\iff\;
  \forall\, p : \mathrm{LPos}(t),\;
  \mathrm{localOk}(\theta,\, t.\mathrm{sub}(p)).
\]
\end{defn}

\paragraph{Lean.}
\begin{lstlisting}[style=lean]
def Ok (theta : T) (t : LTree T) : Prop :=
  forall p : LPos t, (t.sub p).localOk theta
\end{lstlisting}

\paragraph{Haskell.}
\begin{lstlisting}[style=haskell]
ok :: (Ord t) => t -> LTree t -> Bool
ok theta t = all (\p -> localOk theta (sub t p)) (allPositions t)
  where allPositions = ...  -- enumerate all positions recursively
\end{lstlisting}

%===================================================================
\section{Duplicate-stability}
%===================================================================

\subsection{Theorem: \lean{LTree.ok\_sub}}

\begin{thm}[Duplicate-stability of the trust boundary]
If the global boundary rule holds for $t$ at threshold $\theta$, then it
also holds for every subtree:
\[
  \mathrm{Ok}(\theta, t) \;\Rightarrow\;
  \forall\, p : \mathrm{LPos}(t),\;
  \mathrm{Ok}(\theta,\, t.\mathrm{sub}(p)).
\]
The proof is a one-line transport along D4: a position $q$ in
$t.\mathrm{sub}(p)$ corresponds to the position $t.\mathrm{shift}(p, q)$ in
$t$; D4 identifies the subtrees; so the check at $(t, t.\mathrm{shift}(p,
q))$ \emph{is} the check at $(t.\mathrm{sub}(p), q)$.
\end{thm}

\paragraph{Lean.}
\begin{lstlisting}[style=lean]
theorem ok_sub (theta : T) (t : LTree T) (h : Ok theta t)
    (p : LPos t) : Ok theta (t.sub p) :=
  fun q => t.sub_shift p q => h (t.shift p q)
\end{lstlisting}

\paragraph{Haskell.}
\begin{lstlisting}[style=haskell]
-- ok_sub: ok theta t -> ok theta (sub t p)
--
-- For any position q in (sub t p):
--   localOk theta (sub (sub t p) q)
--   = localOk theta (sub t (shift t p q))    -- by D4
--   = h (shift t p q)                        -- from Ok theta t
\end{lstlisting}

\begin{fidelitynote}
In Lean, the transport $\lean{t.sub\_shift\,p\,q} \blacktriangleright
\lean{h\,(t.shift\,p\,q)}$ coerces the hypothesis along the D4 equality.
In Haskell (or informal mathematics) this coercion is invisible because the
identity of the tree doesn't depend on its type index.
\end{fidelitynote}


\subsection{Theorem: \lean{LTree.ok\_iff\_all\_sub}}

\begin{thm}[Ok iff all sub-trees are Ok]
\[
  \mathrm{Ok}(\theta, t) \;\iff\;
  \forall\, p : \mathrm{LPos}(t),\;
  \mathrm{Ok}(\theta,\, t.\mathrm{sub}(p)).
\]
The forward direction is \lean{ok\_sub}; the backward direction reads the
root component (\lean{here}).
\end{thm}

\paragraph{Lean.}
\begin{lstlisting}[style=lean]
theorem ok_iff_all_sub (theta : T) (t : LTree T) :
    Ok theta t <-> forall p : LPos t, Ok theta (t.sub p)
\end{lstlisting}

\paragraph{Haskell.}
\begin{lstlisting}[style=haskell]
-- ok theta t <=> all (\p -> ok theta (sub t p)) (allPositions t)
-- Forward: ok_sub; Backward: apply at Here (sub t Here = t).
\end{lstlisting}

%===================================================================
\section{The deployed trust chain}
%===================================================================

\subsection{Inductive type: \lean{Trust.TrustLevel}}

\begin{defn}[Deployed trust chain]
The seven-level chain used in production registries (schema v3/v4):
\[
  \mathtt{unchecked} < \mathtt{speculative} < \mathtt{computed} <
  \mathtt{proved} < \mathtt{peerReviewed} < \mathtt{published} <
  \mathtt{leanVerified}.
\]
\lean{unchecked} is $\bot$, modelling both \texttt{in-progress} and
\texttt{unclassified} (the validator's CAP\_RANK $-1$).  \texttt{dead-end}
is the \lean{dead} flag, not a chain element.
\end{defn}

\paragraph{Lean.}
\begin{lstlisting}[style=lean]
inductive TrustLevel where
  | unchecked
  | speculative
  | computed
  | proved
  | peerReviewed
  | published
  | leanVerified
deriving DecidableEq
\end{lstlisting}

\paragraph{Haskell.}
\begin{lstlisting}[style=haskell]
data TrustLevel
  = Unchecked | Speculative | Computed | Proved
  | PeerReviewed | Published | LeanVerified
  deriving (Eq, Ord, Show)
\end{lstlisting}


\subsection{Definition: \lean{TrustLevel.rank} and \lean{instance : Chain TrustLevel}}

\begin{defn}[Numeric rank and chain instance]
The rank assigns $\mathtt{unchecked} \mapsto 0,\ldots,\mathtt{leanVerified}
\mapsto 6$; the order $a \le b$ is $\mathrm{rank}(a) \le \mathrm{rank}(b)$.
The \lean{Chain TrustLevel} instance makes the certificate lemma, attainment,
and lax transport from \lean{Trust.lean} available for real registry labels.
\end{defn}

\paragraph{Lean.}
\begin{lstlisting}[style=lean]
def TrustLevel.rank : TrustLevel -> Nat
  | .unchecked    => 0 | .speculative => 1
  | .computed     => 2 | .proved      => 3
  | .peerReviewed => 4 | .published   => 5
  | .leanVerified => 6

instance : Chain TrustLevel where
  le a b := a.rank <= b.rank
  bot := .unchecked
  -- le_refl, le_trans, le_antisymm, le_total, bot_le, dec all proved
\end{lstlisting}

\paragraph{Haskell.}
\begin{lstlisting}[style=haskell]
-- Derived Ord matches the rank ordering.
-- bot = Unchecked.
\end{lstlisting}


\subsection{Abbreviation: \lean{Trust.RegistryOk} and theorem \lean{registryOk\_sub}}

\begin{defn}[Deployed boundary rule]
The \textbf{registry boundary rule} at the deployed threshold \lean{proved}:
\[
  \mathrm{RegistryOk}(t) \;\iff\; \mathrm{Ok}(\mathtt{proved},\, t).
\]
\end{defn}

\begin{thm}[Deployed duplicate-stability]
If $\mathrm{RegistryOk}(t)$ and $p : \mathrm{LPos}(t)$, then
$\mathrm{RegistryOk}(t.\mathrm{sub}(p))$.
\end{thm}

\paragraph{Lean.}
\begin{lstlisting}[style=lean]
abbrev RegistryOk (t : LTree TrustLevel) : Prop :=
  t.Ok TrustLevel.proved

theorem registryOk_sub (t : LTree TrustLevel) (h : RegistryOk t)
    (p : LTree.LPos t) : RegistryOk (t.sub p) :=
  LTree.ok_sub TrustLevel.proved t h p
\end{lstlisting}

\paragraph{Haskell.}
\begin{lstlisting}[style=haskell]
registryOk :: LTree TrustLevel -> Bool
registryOk t = ok Proved t

-- registryOk_sub: registryOk t => registryOk (sub t p)
\end{lstlisting}

%===================================================================
\section{The loop experiment's trust chain}
%===================================================================

\subsection{Inductive type: \lean{LoopTrust.LoopTrustLevel}}

\begin{defn}[Loop trust chain]
The six-level chain for the Warnaar positivity loop experiment:
\[
  \mathtt{unchecked} < \mathtt{speculative} < \mathtt{computed} <
  \mathtt{proved} < \mathtt{verified} < \mathtt{leanVerified}.
\]
Differs from \lean{TrustLevel} in that \lean{verified} replaces
\lean{peerReviewed} and \lean{published}: a \lean{verified} node means an
independent verifier (or Robin) hostile-reviewed the proof.
\end{defn}

\paragraph{Lean.}
\begin{lstlisting}[style=lean]
inductive LoopTrustLevel where
  | unchecked | speculative | computed
  | proved    | verified    | leanVerified
deriving DecidableEq
\end{lstlisting}

\paragraph{Haskell.}
\begin{lstlisting}[style=haskell]
data LoopTrustLevel
  = LUnchecked | LSpeculative | LComputed
  | LProved    | LVerified    | LLeanVerified
  deriving (Eq, Ord, Show)
\end{lstlisting}


\subsection{Definition: \lean{LoopTrustLevel.rank} and \lean{instance : Chain LoopTrustLevel}}

\begin{defn}[Numeric rank and chain instance for loop trust]
Rank runs from 0 (\lean{unchecked}) to 5 (\lean{leanVerified}).  The
\lean{Chain LoopTrustLevel} instance brings the full trust-system theory to
bear on loop registries.
\end{defn}

\paragraph{Lean.}
\begin{lstlisting}[style=lean]
def LoopTrustLevel.rank : LoopTrustLevel -> Nat
  | .unchecked    => 0 | .speculative  => 1
  | .computed     => 2 | .proved       => 3
  | .verified     => 4 | .leanVerified => 5

instance : Chain LoopTrustLevel where
  le a b := a.rank <= b.rank
  bot := .unchecked
  -- le_refl, le_trans, le_antisymm, le_total, bot_le, dec all proved
\end{lstlisting}

\paragraph{Haskell.}
\begin{lstlisting}[style=haskell]
-- Derived Ord matches rank ordering; bot = LUnchecked.
\end{lstlisting}


\subsection{Abbreviation: \lean{LoopRegistryOk} and theorem \lean{loopRegistryOk\_sub}}

\begin{defn}[Loop boundary rule]
\[
  \mathrm{LoopRegistryOk}(t) \;\iff\; \mathrm{Ok}(\mathtt{proved},\, t)
\]
for $t : \lean{LTree}\,\lean{LoopTrustLevel}$.
\end{defn}

\begin{thm}[Loop duplicate-stability]
$\mathrm{LoopRegistryOk}(t)$ implies $\mathrm{LoopRegistryOk}(t.\mathrm{sub}(p))$
for every $p$.
\end{thm}

\paragraph{Lean.}
\begin{lstlisting}[style=lean]
abbrev LoopRegistryOk (t : Containers.LTree LoopTrustLevel) : Prop :=
  t.Ok LoopTrustLevel.proved

theorem loopRegistryOk_sub (t : Containers.LTree LoopTrustLevel)
    (h : LoopRegistryOk t) (p : Containers.LTree.LPos t) :
    LoopRegistryOk (t.sub p) :=
  Containers.LTree.ok_sub LoopTrustLevel.proved t h p
\end{lstlisting}

\paragraph{Haskell.}
\begin{lstlisting}[style=haskell]
loopRegistryOk :: LTree LoopTrustLevel -> Bool
loopRegistryOk t = ok LProved t

-- loopRegistryOk_sub: loopRegistryOk t => loopRegistryOk (sub t p)
\end{lstlisting}

%===================================================================
\section{Translation conventions}
%===================================================================

\begin{enumerate}
\item \textbf{Dependent type families.}  Positions \lean{LPos t} depend on
  the tree $t$ in Lean; in Haskell all trees have the same position type,
  losing this dependency.

\item \textbf{Transport.}  D5 (\lean{shift\_assoc}) uses transport along the
  D4 equality.  The listings approximate \lean{|>} (the Lean transport operator) by \texttt{=>} (meaning
  ``cast via this proof''); no Haskell counterpart exists.

\item \textbf{Duplicate-stability via D4.}  The one-line proof of
  \lean{ok\_sub} is $q \mapsto \lean{sub\_shift\,p\,q} \blacktriangleright
  \lean{h\,(shift\,p\,q)}$: a D4 transport applied to the hypothesis.
  In Haskell, D4 holds computationally and no coercion is needed.

\item \textbf{Proof bodies omitted.}  Lean checks all proofs.

\item \textbf{Namespace structure.}  \lean{TrustLabel}, \lean{LTree},
  \lean{trustRegistry} are in \lean{Containers}; \lean{TrustLevel} and
  \lean{RegistryOk} are in \lean{Containers.Trust};
  \lean{LoopTrustLevel} and \lean{LoopRegistryOk} are in
  \lean{Containers.LoopTrust}.
\end{enumerate}

\end{document}
